\documentclass[11pt,a4paper]{article}
\usepackage[utf8]{inputenc}
\usepackage[T1]{fontenc}
\usepackage[ngerman]{babel}
\usepackage[margin=2.5cm]{geometry}
\usepackage{hyperref}
\usepackage{xcolor}
\usepackage{listings}
\usepackage{enumitem}
\usepackage{booktabs}
\usepackage{longtable}
\usepackage{amsmath,amssymb}
\usepackage{fancyhdr}
\usepackage{titlesec}
\definecolor{isls-blue}{HTML}{1a5276}
\definecolor{isls-orange}{HTML}{e67e22}
\definecolor{isls-green}{HTML}{27ae60}
\definecolor{isls-red}{HTML}{c0392b}
\definecolor{isls-gray}{HTML}{7f8c8d}
\definecolor{code-bg}{HTML}{f7f9fb}
\lstdefinestyle{rust}{language=C,basicstyle=\ttfamily\small,keywordstyle=\color{isls-blue}\bfseries,commentstyle=\color{isls-gray}\itshape,stringstyle=\color{isls-green},backgroundcolor=\color{code-bg},frame=single,rulecolor=\color{isls-gray!30},breaklines=true,showstringspaces=false,tabsize=4,morekeywords={fn,let,mut,pub,struct,impl,enum,use,mod,self,Self,crate,super,where,trait,for,in,if,else,match,return,async,await,Ok,Err,Some,None,Vec,String,HashMap,BTreeMap,BTreeSet,Result,Option,PathBuf,f64,f32,usize,bool,u32,u64,i64,Box,dyn,Arc}}
\lstdefinestyle{bash}{language=bash,basicstyle=\ttfamily\small,backgroundcolor=\color{code-bg},frame=single,rulecolor=\color{isls-gray!30},breaklines=true,showstringspaces=false}
\pagestyle{fancy}\fancyhf{}
\fancyhead[L]{\small\textcolor{isls-gray}{ISLS N1 Spec --- Normsystem-Kern}}
\fancyhead[R]{\small\textcolor{isls-gray}{v1.0.0 --- \today}}
\fancyfoot[C]{\thepage}
\titleformat{\section}{\Large\bfseries\color{isls-blue}}{\thesection}{1em}{}
\titleformat{\subsection}{\large\bfseries\color{isls-blue!80}}{\thesubsection}{1em}{}
\titleformat{\subsubsection}{\normalsize\bfseries\color{isls-blue!60}}{\thesubsubsection}{1em}{}
\hypersetup{colorlinks=true,linkcolor=isls-blue,urlcolor=isls-orange}
\begin{document}
\begin{titlepage}\centering\vspace*{2.5cm}
{\Huge\bfseries\color{isls-blue} ISLS N1 Specification}\\[0.5cm]
{\LARGE\color{isls-orange} Normsystem-Kern}\\[0.3cm]
{\large\color{isls-gray} Relationen, Valenz, Gruppen, Konfigurationsenergie,\\Optimale Komposition, Studio-Visualisierung}\\[2cm]
{\large\begin{tabular}{rl}
\textbf{System:} & ISLS --- Intelligent Semantic Ledger Substrate \\
\textbf{Version:} & v7.1 (N1) \\
\textbf{Author:} & Sebastian Klemm \\
\textbf{Date:} & \today \\
\textbf{Status:} & Specification \\
\textbf{Depends:} & I5 (complete), F1 (complete) \\
\textbf{Branch:} & \texttt{claude/implement-isls-n1-spec} \\
\textbf{Basis:} & Das Universelle Normsystem (Formalisierung) \\
\textbf{Regel 13:} & Alles im Studio sichtbar.
\end{tabular}}
\vfill{\small\color{isls-gray}Mithras Substructure University\\License: MIT}
\end{titlepage}
\tableofcontents\newpage

% ======================================================================
\section{Executive Summary}
% ======================================================================

Die Norm-Bibliothek ist aktuell eine flache Liste. N1 installiert die algebraische Struktur darauf: Relationen (Kompatibilit\"at, Konflikt, Abh\"angigkeit), Valenz, Norm-Gruppen, Konfigurationsenergie und einen Greedy-Optimierer der f\"ur eine Anforderung die stabilste Konfiguration findet.

Sechs Work Packages:

\begin{enumerate}[nosep]
\item \textbf{W1: Relationen-Engine} --- Berechnung von $\sim$, $\perp$, $\to$ aus Resonites.
\item \textbf{W2: Valenz + Gruppen} --- Bindungsf\"ahigkeit und Periodensystem.
\item \textbf{W3: Konfigurationsenergie} --- $E(\mathcal{K})$ berechnen.
\item \textbf{W4: Greedy-Optimierer} --- Optimale Konfiguration f\"ur Anforderung.
\item \textbf{W5: Norm-Selektion im Forge} --- Forge nutzt den Optimierer statt flache Aktivierung.
\item \textbf{W6: Studio-Visualisierung} --- Periodensystem, Konfigurationsenergie, Optimierer im GUI.
\end{enumerate}

% ======================================================================
\section{W1: Relationen-Engine}
\label{sec:w1}
% ======================================================================

\subsection{Datenstrukturen}

\begin{lstlisting}[style=rust]
/// Relations between two norms.
#[derive(Debug, Clone, Serialize, Deserialize)]
pub struct NormRelation {
    pub norm_a: String,
    pub norm_b: String,
    pub relation: RelationType,
    pub strength: f64,        // 0.0-1.0, how strong the relation
    pub shared_resonites: Vec<String>,  // resonite names in common
}

#[derive(Debug, Clone, Copy, PartialEq, Eq, Serialize, Deserialize)]
pub enum RelationType {
    Compatible,     // can coexist, may benefit each other
    Dependent,      // norm_a requires norm_b
    Conflicting,    // cannot coexist
    Independent,    // no interaction
}

/// Complete relation matrix for all norms.
#[derive(Debug, Clone, Serialize, Deserialize)]
pub struct NormRelationMatrix {
    pub relations: Vec<NormRelation>,
    pub computed_at: String,
    pub norm_count: usize,
}
\end{lstlisting}

\subsection{Berechnung der Relationen}

\begin{lstlisting}[style=rust]
pub fn compute_relations(norms: &[Norm]) -> NormRelationMatrix {
    let mut relations = Vec::new();
    
    for i in 0..norms.len() {
        for j in (i+1)..norms.len() {
            let rel = compute_pair_relation(&norms[i], &norms[j]);
            if rel.relation != RelationType::Independent {
                relations.push(rel);
            }
        }
    }
    
    NormRelationMatrix {
        relations,
        computed_at: now_iso(),
        norm_count: norms.len(),
    }
}

fn compute_pair_relation(a: &Norm, b: &Norm) -> NormRelation {
    let a_exports = extract_exports(a);  // Fn names, Type names
    let b_exports = extract_exports(b);
    let a_imports = extract_imports(a);   // required types/fns
    let b_imports = extract_imports(b);
    
    // Dependency: a imports what b exports
    let a_needs_b: Vec<String> = a_imports.intersection(&b_exports)
        .cloned().collect();
    let b_needs_a: Vec<String> = b_imports.intersection(&a_exports)
        .cloned().collect();
    
    // Conflict: both export same name with different signatures
    let shared_exports: Vec<String> = a_exports.intersection(&b_exports)
        .cloned().collect();
    let conflicting = shared_exports.iter().any(|name| {
        let sig_a = get_signature(a, name);
        let sig_b = get_signature(b, name);
        sig_a != sig_b  // same name, different signature = conflict
    });
    
    // Compatible: shared resonite classes
    let a_classes = classify_norm_resonites(a);
    let b_classes = classify_norm_resonites(b);
    let shared_classes: Vec<String> = a_classes.intersection(&b_classes)
        .cloned().collect();
    
    let (relation, strength, shared) = if conflicting {
        (RelationType::Conflicting, 
         shared_exports.len() as f64 / a_exports.len().max(1) as f64,
         shared_exports)
    } else if !a_needs_b.is_empty() {
        (RelationType::Dependent, 
         a_needs_b.len() as f64 / a_imports.len().max(1) as f64,
         a_needs_b)
    } else if !b_needs_a.is_empty() {
        // Reverse dependency: b depends on a
        // Store as a separate relation with norm_a=b, norm_b=a
        (RelationType::Dependent, 
         b_needs_a.len() as f64 / b_imports.len().max(1) as f64,
         b_needs_a)
    } else if !shared_classes.is_empty() {
        (RelationType::Compatible, 
         shared_classes.len() as f64 / 
            a_classes.union(&b_classes).count().max(1) as f64,
         shared_classes)
    } else {
        (RelationType::Independent, 0.0, vec![])
    };
    
    NormRelation {
        norm_a: a.id.clone(),
        norm_b: b.id.clone(),
        relation,
        strength,
        shared_resonites: shared,
    }
}
\end{lstlisting}

\subsection{Import/Export-Extraktion aus Normen}

Die bestehenden Norm-Strukturen haben \texttt{layers} und \texttt{expected\_resonites}. Daraus:

\begin{lstlisting}[style=rust]
fn extract_exports(norm: &Norm) -> BTreeSet<String> {
    // Alle Fn-Resonites und Type-Resonites die die Norm definiert
    norm.resonites.iter()
        .filter(|r| matches!(r.kind, ResoniteKind::Fn | ResoniteKind::Type))
        .map(|r| r.name.clone())
        .collect()
}

fn extract_imports(norm: &Norm) -> BTreeSet<String> {
    // Alle Import-Resonites + Typen die in Fn-Signaturen
    // referenziert aber nicht in der Norm definiert werden
    let defined: BTreeSet<String> = extract_exports(norm);
    norm.resonites.iter()
        .filter(|r| matches!(r.kind, ResoniteKind::Import))
        .map(|r| r.name.clone())
        .chain(
            norm.resonites.iter()
                .filter_map(|r| r.return_type.as_ref())
                .filter(|t| !defined.contains(*t))
                .cloned()
        )
        .collect()
}
\end{lstlisting}

\subsection{Persistenz}

Die Relationsmatrix wird in \texttt{\textasciitilde/.isls/relations.json} gespeichert. Neuberechnung bei:
\begin{itemize}[nosep]
\item Neue Norm promoviert (auto oder injiziert)
\item \texttt{isls norms relations --recompute}
\item Alle 100 neuen Beobachtungen (automatisch)
\end{itemize}

\subsection{CLI}

\begin{lstlisting}[style=bash]
isls norms relations

ISLS Norm Relations (43 norms, 127 relations)
---
Compatible:  89
Dependent:   31
Conflicting:  7

Top Valenz:
  ISLS-NORM-0042 (CrudEntity)       V=18
  ISLS-NORM-0088 (JWT-Auth)         V=12
  ISLS-NORM-INFRA-DB                V=15
  ISLS-NORM-AUTO-0003               V=8

Konflikte:
  JWT-Auth <> Session-Auth          (shared: authenticate)
  Diesel-ORM <> SeaORM              (shared: query, model)
  REST-API <> GraphQL-API           (shared: endpoint)

Abhaengigkeiten:
  CrudEntity -> Database
  JWT-Auth -> CrudEntity (User)
  Pagination -> CrudEntity
  Search -> CrudEntity
  FileUpload -> CrudEntity
\end{lstlisting}

\subsection{API Endpoint}

\begin{lstlisting}[style=bash]
GET /api/norms/relations
Response: NormRelationMatrix as JSON

GET /api/norms/relations?norm=ISLS-NORM-0042
Response: Only relations involving this norm
\end{lstlisting}

% ======================================================================
\section{W2: Valenz und Norm-Gruppen}
\label{sec:w2}
% ======================================================================

\subsection{Valenz-Berechnung}

\begin{lstlisting}[style=rust]
pub fn compute_valence(
    norm_id: &str,
    relations: &NormRelationMatrix
) -> usize {
    relations.relations.iter()
        .filter(|r| 
            (r.norm_a == norm_id || r.norm_b == norm_id) &&
            r.relation == RelationType::Compatible
        )
        .count()
}
\end{lstlisting}

\subsection{Norm-Gruppen-Erkennung}

\begin{lstlisting}[style=rust]
/// A group of exchangeable norms (like isotopes).
#[derive(Debug, Clone, Serialize, Deserialize)]
pub struct NormGroup {
    pub id: String,           // "GROUP-AUTH", "GROUP-ORM", etc.
    pub name: String,         // "Authentication", "ORM", etc.
    pub members: Vec<String>, // norm IDs
    pub shared_deps: Vec<String>,  // dependencies all members share
    pub role: String,         // what function the group fulfills
}

pub fn detect_groups(
    norms: &[Norm],
    relations: &NormRelationMatrix
) -> Vec<NormGroup> {
    let mut groups = Vec::new();
    let mut assigned: BTreeSet<String> = BTreeSet::new();
    
    // Find conflict cliques: norms that conflict with each other
    // are alternatives for the same role
    for i in 0..norms.len() {
        if assigned.contains(&norms[i].id) { continue; }
        
        let mut group_members = vec![norms[i].id.clone()];
        
        for j in (i+1)..norms.len() {
            if assigned.contains(&norms[j].id) { continue; }
            
            // Check: does j conflict with all current members?
            let conflicts_with_all = group_members.iter().all(|m| {
                relations.relations.iter().any(|r| {
                    r.relation == RelationType::Conflicting &&
                    ((r.norm_a == *m && r.norm_b == norms[j].id) ||
                     (r.norm_b == *m && r.norm_a == norms[j].id))
                })
            });
            
            // AND: shares the same dependencies
            let same_deps = get_dependencies(&norms[i].id, relations) ==
                           get_dependencies(&norms[j].id, relations);
            
            if conflicts_with_all && same_deps {
                group_members.push(norms[j].id.clone());
            }
        }
        
        if group_members.len() > 1 {
            let shared_deps = get_dependencies(
                &group_members[0], relations
            );
            let name = infer_group_name(&group_members, norms);
            
            for m in &group_members {
                assigned.insert(m.clone());
            }
            
            groups.push(NormGroup {
                id: format!("GROUP-{}", groups.len() + 1),
                name,
                members: group_members,
                shared_deps,
                role: String::new(), // filled by classification
            });
        }
    }
    
    groups
}
\end{lstlisting}

\subsection{Persistenz}

\texttt{\textasciitilde/.isls/groups.json}. Neuberechnung zusammen mit der Relationsmatrix.

\subsection{CLI}

\begin{lstlisting}[style=bash]
isls norms groups

ISLS Norm Groups (4 groups, 12 grouped norms, 31 singletons)
---
GROUP-1  "Authentication"   [JWT-Auth, Session-Auth]
  Role: User authentication
  Shared deps: Database, User-Model
  
GROUP-2  "ORM"              [Diesel, SeaORM, SQLx-raw]
  Role: Database abstraction
  Shared deps: Database-Connection
  
GROUP-3  "API-Style"        [REST, GraphQL]
  Role: External interface
  Shared deps: CrudEntity, Auth
  
GROUP-4  "Cache"            [Redis, InMemory]
  Role: Data caching
  Shared deps: Configuration
\end{lstlisting}

% ======================================================================
\section{W3: Konfigurationsenergie}
\label{sec:w3}
% ======================================================================

\begin{lstlisting}[style=rust]
/// Energy of a configuration (set of norms).
pub fn configuration_energy(
    config: &[String],    // norm IDs
    relations: &NormRelationMatrix,
    fitness: &HashMap<String, f64>,
) -> ConfigEnergy {
    let mut conflict_energy = 0.0;
    let mut deficit_energy = 0.0;
    let mut fitness_entropy = 0.0;
    let mut conflicts = Vec::new();
    let mut missing_deps = Vec::new();
    
    // 1. Conflict energy
    for r in &relations.relations {
        if r.relation == RelationType::Conflicting {
            let a_in = config.contains(&r.norm_a);
            let b_in = config.contains(&r.norm_b);
            if a_in && b_in {
                conflict_energy += r.strength;
                conflicts.push((r.norm_a.clone(), r.norm_b.clone()));
            }
        }
    }
    
    // 2. Deficit energy (missing dependencies)
    for r in &relations.relations {
        if r.relation == RelationType::Dependent {
            let a_in = config.contains(&r.norm_a);
            let b_in = config.contains(&r.norm_b);
            if a_in && !b_in {
                deficit_energy += r.strength;
                missing_deps.push((r.norm_a.clone(), r.norm_b.clone()));
            }
        }
    }
    
    // 3. Fitness entropy
    for norm_id in config {
        let phi = fitness.get(norm_id).copied().unwrap_or(0.5);
        fitness_entropy += 1.0 - phi;
    }
    
    ConfigEnergy {
        total: conflict_energy * 10.0 + deficit_energy * 5.0 
               + fitness_entropy,
        conflict: conflict_energy,
        deficit: deficit_energy,
        fitness_entropy,
        conflicts,
        missing_deps,
        is_valid: conflicts.is_empty() && missing_deps.is_empty(),
    }
}

#[derive(Debug, Clone, Serialize, Deserialize)]
pub struct ConfigEnergy {
    pub total: f64,
    pub conflict: f64,
    pub deficit: f64,
    pub fitness_entropy: f64,
    pub conflicts: Vec<(String, String)>,
    pub missing_deps: Vec<(String, String)>,
    pub is_valid: bool,
}
\end{lstlisting}

\subsection{CLI}

\begin{lstlisting}[style=bash]
isls norms energy --norms ISLS-NORM-0042,ISLS-NORM-0088,ISLS-NORM-0096

Configuration Energy Analysis
---
Norms: CrudEntity, JWT-Auth, Pagination
Energy: 1.35
  Conflict:        0.00  (keine Konflikte)
  Deficit:         0.85  (fehlend: Database, ErrorHandling)
  Fitness-Entropy: 0.50

Status: UNVOLLSTAENDIG
  Fehlende Abhaengigkeiten:
    CrudEntity -> Database
    JWT-Auth -> Database

Empfehlung: +Database, +ErrorHandling -> Energy: 0.30
\end{lstlisting}

% ======================================================================
\section{W4: Greedy-Optimierer}
\label{sec:w4}
% ======================================================================

\begin{lstlisting}[style=rust]
/// Find optimal norm configuration for a set of requirements.
pub fn optimize_configuration(
    requirements: &[ResoniteClass],  // what the app needs
    norms: &[Norm],
    relations: &NormRelationMatrix,
    fitness: &HashMap<String, f64>,
    groups: &[NormGroup],
) -> OptimalConfig {
    let mut selected: Vec<String> = Vec::new();
    let mut covered: BTreeSet<ResoniteClass> = BTreeSet::new();
    let mut remaining = BTreeSet::from_iter(requirements.iter().cloned());
    
    // Sort norms by fitness (descending)
    let mut candidates: Vec<(&Norm, f64)> = norms.iter()
        .map(|n| (n, *fitness.get(&n.id).unwrap_or(&0.5)))
        .collect();
    candidates.sort_by(|a, b| b.1.partial_cmp(&a.1).unwrap());
    
    while !remaining.is_empty() {
        let mut best: Option<(&Norm, usize)> = None;
        
        for (norm, _phi) in &candidates {
            if selected.contains(&norm.id) { continue; }
            
            // Check: does this norm conflict with any selected?
            let has_conflict = selected.iter().any(|s| {
                relations.relations.iter().any(|r| {
                    r.relation == RelationType::Conflicting &&
                    ((r.norm_a == *s && r.norm_b == norm.id) ||
                     (r.norm_b == *s && r.norm_a == norm.id))
                })
            });
            if has_conflict { continue; }
            
            // How many remaining requirements does it cover?
            let norm_classes = classify_norm_resonites(norm);
            let covers: usize = remaining.intersection(&norm_classes)
                .count();
            
            if covers > 0 {
                if best.is_none() || covers > best.unwrap().1 {
                    best = Some((norm, covers));
                }
            }
        }
        
        match best {
            Some((norm, _)) => {
                // Add norm + its dependency closure
                selected.push(norm.id.clone());
                let norm_classes = classify_norm_resonites(norm);
                for c in &norm_classes {
                    remaining.remove(c);
                    covered.insert(c.clone());
                }
                
                // Add dependencies
                let deps = transitive_closure(&norm.id, relations);
                for dep in deps {
                    if !selected.contains(&dep) {
                        selected.push(dep);
                    }
                }
            }
            None => break, // no norm covers remaining requirements
        }
    }
    
    // For each group: if multiple members selected, keep highest fitness
    resolve_group_conflicts(&mut selected, groups, fitness);
    
    let energy = configuration_energy(&selected, relations, fitness);
    
    OptimalConfig {
        norms: selected,
        covered: covered.into_iter().collect(),
        uncovered: remaining.into_iter().collect(),
        energy,
    }
}

#[derive(Debug, Clone, Serialize, Deserialize)]
pub struct OptimalConfig {
    pub norms: Vec<String>,
    pub covered: Vec<ResoniteClass>,
    pub uncovered: Vec<ResoniteClass>,
    pub energy: ConfigEnergy,
}
\end{lstlisting}

\subsection{CLI}

\begin{lstlisting}[style=bash]
isls norms optimize --require CrudEntity,Authentication,Pagination,WebSocket

Optimal Configuration
---
Selected (7 norms, Energy: 0.42):
  ISLS-NORM-0042  CrudEntity        phi=0.95
  ISLS-NORM-0088  JWT-Auth          phi=0.93  (Group: Authentication)
  ISLS-NORM-0096  Pagination        phi=0.92
  ISLS-NORM-INFRA-DB  Database      phi=0.90  (dependency)
  ISLS-NORM-INFRA-AUTH ErrorSystem   phi=0.88  (dependency)
  ISLS-NORM-AUTO-0005  WebSocket    phi=0.78
  ISLS-NORM-AUTO-0008  EventBus     phi=0.72  (dependency of WebSocket)

Energy Breakdown:
  Conflict:        0.00
  Deficit:         0.00  (alle Abhaengigkeiten erfuellt)
  Fitness-Entropy: 0.42

Alternatives (Group choices):
  Authentication: JWT-Auth [SELECTED] | Session-Auth (phi=0.85)
  
Uncovered: (none)
\end{lstlisting}

% ======================================================================
\section{W5: Norm-Selektion im Forge}
\label{sec:w5}
% ======================================================================

\subsection{Problem}

Aktuell aktiviert der Forge Normen durch Keyword-Matching: wenn die App-Beschreibung ``authentication'' enth\"alt, werden Auth-Normen aktiviert. Das ist flach --- es ignoriert Abh\"angigkeiten, Konflikte und Energie.

\subsection{L\"osung: Optimierer-basierte Aktivierung}

Der Forge-Prozess bekommt einen neuen Schritt zwischen S0 (Ingest) und S1 (Canon):

\begin{lstlisting}[style=rust]
// In staged_closure.rs, after TOML parsing:

// 1. Extract requirements from entities
let requirements = extract_requirements_from_entities(&entities);
// e.g. [CrudEntity, Authentication, Pagination]

// 2. Run optimizer
let optimal = optimize_configuration(
    &requirements, &norms, &relations, &fitness, &groups
);

// 3. Use optimal.norms as the activated norms
// (instead of keyword-matching)
let activated_norms = optimal.norms;

// 4. Log energy
tracing::info!(
    "Norm selection: {} norms, E={:.2}, {} uncovered",
    activated_norms.len(), optimal.energy.total, 
    optimal.uncovered.len()
);
\end{lstlisting}

\subsection{Anforderungs-Extraktion}

\begin{lstlisting}[style=rust]
fn extract_requirements_from_entities(
    entities: &[EntityDef]
) -> Vec<ResoniteClass> {
    let mut reqs = vec![ResoniteClass::CrudEntity]; // always
    
    // Auth: always for web apps
    reqs.push(ResoniteClass::Authentication);
    
    // Pagination: if any entity might have many entries
    if entities.len() > 1 {
        reqs.push(ResoniteClass::Pagination);
    }
    
    // Check entity fields for hints
    for entity in entities {
        for field in &entity.fields {
            let name = field.name.to_lowercase();
            if name.contains("status") || name.contains("state") {
                reqs.push(ResoniteClass::StateMachine);
            }
            if name.contains("file") || name.contains("image") 
               || name.contains("attachment") {
                reqs.push(ResoniteClass::FileUpload);
            }
            if name.contains("email") {
                reqs.push(ResoniteClass::Notification);
            }
        }
    }
    
    reqs.sort();
    reqs.dedup();
    reqs
}
\end{lstlisting}

% ======================================================================
\section{W6: Studio-Visualisierung}
\label{sec:w6}
% ======================================================================

Regel 13: Alles im Studio sichtbar. Drei neue Visualisierungen:

\subsection{6a: Periodensystem im Entdecken-Modus}

Trigger: ``periodensystem'' oder ``relations'' im Chat.

Darstellung: Normen als Kreise, angeordnet nach Gruppen. Kompatible Verbindungen als gr\"une Linien. Konflikte als rote Linien. Abh\"angigkeiten als Pfeile. Kreisgr\"o{\ss}e $\propto$ Valenz. Farbe $=$ Fitness (gr\"un $>$ 0.8, orange 0.5--0.8, rot $<$ 0.5).

SVG-basiert, inline. Kein D3. Force-Layout nicht n\"otig --- einfaches Grid nach Gruppen sortiert.

\begin{lstlisting}[style=bash]
  Norm-Periodensystem (43 Normen, 4 Gruppen)
  
  [Auth-Gruppe]         [ORM-Gruppe]        [Singletons]
  +------+------+      +------+------+     +------+
  | JWT  | Sess |      | Dsl  | Sea  |     | CRUD |<--+
  | 0.93 | 0.85 |      | 0.88 | 0.91 |     | 0.95 |   |
  +------+------+      +------+------+     +------+   |
        |                    |                  |      |
        +-------- dep -------+---- dep ---------+      |
                                                       |
  [API-Gruppe]           [Cache-Gruppe]     +------+   |
  +------+------+       +------+------+    | Pagi | --+
  | REST | GQL  |       | Rds  | Mem  |    | 0.92 |
  | 0.95 | 0.78 |       | 0.72 | 0.65 |    +------+
  +------+------+       +------+------+
\end{lstlisting}

\subsection{6b: Energie-Rechner im Erschaffen-Modus}

Nach dem Architect-Chat (Entities definiert, vor dem Forge): automatisch die optimale Konfiguration berechnen und im Chat anzeigen:

\begin{lstlisting}[style=bash]
  Optimale Norm-Konfiguration (E=0.42):
    * CrudEntity (0.95)
    * JWT-Auth (0.93)
    * Pagination (0.92)
    * Database (0.90) [auto]
    * ErrorSystem (0.88) [auto]
  
  Alternative: Session-Auth statt JWT (E=0.48)
  
  [Forge]
\end{lstlisting}

Der Nutzer sieht welche Normen aktiviert werden und warum, bevor er Forge dr\"uckt.

\subsection{6c: Konfigurationsenergie im Sensorium}

8. Karte im Sensorium: ``Konfigurationsenergie''. Zeigt f\"ur die letzten 10 Forge-Runs: welche Normen aktiviert, welche Energie, ob Konflikte/Defizite existierten. Liniengraph der Energie \"uber Zeit.

\subsection{Neue Endpoints}

\begin{lstlisting}[style=bash]
GET  /api/norms/relations     -> NormRelationMatrix
GET  /api/norms/groups        -> Vec<NormGroup>
POST /api/norms/energy        -> ConfigEnergy
  Body: { "norms": ["ISLS-NORM-0042", ...] }
POST /api/norms/optimize      -> OptimalConfig
  Body: { "requirements": ["CrudEntity", "Authentication", ...] }
\end{lstlisting}

% ======================================================================
\section{Neue Dateien}
% ======================================================================

\begin{longtable}{lp{10cm}}
\toprule
\textbf{Datei} & \textbf{Inhalt} \\
\midrule
\texttt{isls-norms/src/relations.rs}
  & \texttt{NormRelation}, \texttt{NormRelationMatrix}, \texttt{compute\_relations()}, Import/Export-Extraktion, Persistenz. \\
\texttt{isls-norms/src/groups.rs}
  & \texttt{NormGroup}, \texttt{detect\_groups()}, Gruppenname-Inferenz. \\
\texttt{isls-norms/src/energy.rs}
  & \texttt{ConfigEnergy}, \texttt{configuration\_energy()}, Energie-Berechnung. \\
\texttt{isls-norms/src/optimizer.rs}
  & \texttt{OptimalConfig}, \texttt{optimize\_configuration()}, Greedy-Algorithmus, Anforderungs-Extraktion. \\
\bottomrule
\end{longtable}

% ======================================================================
\section{Modifizierte Dateien}
% ======================================================================

\begin{longtable}{lp{10cm}}
\toprule
\textbf{Datei} & \textbf{\"Anderung} \\
\midrule
\texttt{isls-norms/src/lib.rs}
  & \texttt{pub mod relations; pub mod groups; pub mod energy; pub mod optimizer;} \\
\texttt{isls-forge-llm/src/staged\_closure.rs}
  & Norm-Selektion durch Optimierer statt Keyword-Match. \\
\texttt{isls-cli/src/main.rs}
  & Subcommands: \texttt{isls norms relations}, \texttt{isls norms groups}, \texttt{isls norms energy}, \texttt{isls norms optimize}. \\
\texttt{isls-gateway/src/lib.rs}
  & 4 neue API-Routen. \\
\texttt{isls-gateway/src/discover.rs}
  & Handler f\"ur Relations/Groups/Energy/Optimize. \\
\texttt{isls-gateway/src/static/studio.html}
  & Periodensystem, Energie-Rechner, Sensorium-Karte. \\
\bottomrule
\end{longtable}

% ======================================================================
\section{Constraints for Claude Code}
% ======================================================================

\begin{enumerate}
\item \textbf{Regel 13:} Periodensystem, Energie-Rechner und Konfigurationsenergie-Graph m\"ussen im Studio sichtbar sein.
\item \textbf{Relationen werden berechnet, nicht manuell definiert.} Aus Resonite-\"Uberlappung, nicht aus einer hardcodierten Tabelle.
\item \textbf{Greedy-Optimierer}, kein ILP. Einfach, schnell, gut genug f\"ur $|\mathcal{N}| \leq 200$.
\item \textbf{Persistenz:} \texttt{relations.json}, \texttt{groups.json} neben \texttt{norms.json}. Neuberechnung bei Promotion.
\item \textbf{Forge-Integration:} Der Optimierer ersetzt die Keyword-Aktivierung. Aber Fallback: wenn der Optimierer keine Konfiguration findet (z.B. leere Norm-Bibliothek), verwende die alte Keyword-Logik.
\item \textbf{Energie-Gewichte:} Konflikt $\times 10$, Defizit $\times 5$, Fitness-Entropie $\times 1$. Konflikte wiegen am schwersten.
\item \textbf{SVG-Periodensystem:} Grid-Layout nach Gruppen. Kein Force-Layout. Einfach, deterministisch.
\item \textbf{Do not modify} isls-reader, isls-code-topo, isls-hypercube, isls-types, isls-chat, isls-merkaba.
\item Alle bestehenden Tests m\"ussen passen.
\item Branch: \texttt{claude/implement-isls-n1-spec}.
\item Commit: \texttt{N1/W\{n\}: <description>}.
\end{enumerate}

% ======================================================================
\section{Acceptance Criteria}
% ======================================================================

\begin{longtable}{clp{9cm}}
\toprule
\# & Pass & Kriterium \\
\midrule
1 & & \texttt{cargo build --workspace} kompiliert. \\
2 & & \texttt{cargo test --workspace} besteht. \\
3 & & \texttt{isls norms relations} zeigt Kompatibilit\"aten, Konflikte, Abh\"angigkeiten. \\
4 & & \texttt{isls norms groups} erkennt mindestens 1 Gruppe (wenn Konflikte existieren). \\
5 & & \texttt{isls norms energy --norms A,B,C} berechnet Energie. \\
6 & & \texttt{isls norms optimize --require CrudEntity,Authentication} findet optimale Konfiguration. \\
7 & & Forge nutzt Optimierer f\"ur Norm-Selektion. \\
8 & & Studio: ``periodensystem'' zeigt Norm-Graph. \\
9 & & Studio: Vor Forge wird optimale Konfiguration angezeigt. \\
10 & & Studio: Sensorium hat Konfigurationsenergie-Karte. \\
11 & & Alle bestehenden Tests bestehen. \\
\bottomrule
\end{longtable}

\vfill
\begin{center}
\color{isls-gray}\rule{0.5\textwidth}{0.4pt}\\[0.5em]
{\small Ende der N1-Spezifikation.\\[0.3em]
Nicht die Normen sind das System.\\Die Relationen zwischen ihnen sind es.\\Und jetzt kann die Maschine sie sehen.}
\end{center}
\end{document}
